\uuid{ynIC}
\titre{Convergence}
\auteur{bodin ; hanani ; mzari}
\organisation{exo7}
\licence{CC BY-NC-SA 4.0}
\datecreate{2020-01-20}
\niveau{L1}
\module{Analyse}
\chapitre{Suite}
\sousChapitre{Autre}
\difficulte{1}
\theme{Suites, convergence}
\exo7qcmid{lille-1.suites.010}

\contenu{
\begin{question}
\qclassification{121.00}
\qtags{motcle=convergence}
Soit $\displaystyle u_n=\sqrt[n]{3+\cos n}$. Quelles sont les bonnes réponses ?
\begin{answers}
    \good{La suite $(u_n)$ est bornée.}
    \good{$\displaystyle \lim _{n\to +\infty}u_n=1$}
    \bad{La suite $(u_n)$ est croissante.}
    \bad{La suite $(u_n)$ est divergente.}
\end{answers}
\begin{explanations}
On a : $\displaystyle \sqrt[n]{2}\leq u_n\leq \sqrt[n]{4}$. Or
$$\lim _{n\to +\infty}\sqrt[n]{2}=1=\lim _{n\to +\infty}\sqrt[n]{4}.$$
Donc, le théorème d'encadrement implique que $(u_n)$ converge et que sa limite est $1$.
\end{explanations}
\end{question}
}
